% Part: proof-theory
% Chapter: natural-deduction
% Section: translation-N2i

\documentclass[../../../include/open-logic-section]{subfiles}

\begin{document}

\olfileid{pt}{nat}{tni}

\olsection{\Log{N2}માંથી \Log{G2}માં રૂપાંતર}

\begin{prop}\ollabel{prop:N2-to-G2}
જો $\Log{N2c}\ (\Log{N2i}) \Proves
\Gamma \Sequent !A$,
તો~$\Log{G2c}\ (\Log{G2i}) + \Cut \Proves \Gamma'
\Sequent \Delta$,
જ્યાં~$\Gamma$માંથી
ચિહ્નો દૂર કરતાં
મળતો !!{formula}sનો
બહુગણ~$\Gamma'$ છે,
અને~$!A$ એ~$\lfalse$
ન હોય તો~$\Delta = \{!A\}$,
જ્યારે હોય તો
$\Delta = \emptyset$.
\sourcecorrection{OLMD-249}{મૂળના પરિણામમાં G2c પછી કૌંસમાં N2i લખાયું છે; અનુરૂપ અંતઃપ્રજ્ઞાવાદી સિક્વન્ટ-કલન G2i છે.}
\end{prop}

\begin{proof}
ફરી~$\Gamma \Sequent !A$ની
!!{proof}~$\delta$ની
ઊંચાઈ પર આગમન કરીએ.
જો~$\delta$ એ
$x:!A \Sequent !A$
સ્વયંસિદ્ધ હોય,
તો~$\pi$ એ અનુરૂપ
$!A \Sequent !A$
સ્વયંસિદ્ધ છે;
$!A \ident \lfalse$
હોય ત્યારે
$\lfalse \Sequent \ $
લેવું.

જો~$\delta$નો અંત
અનુમાન-નિયમમાં થાય,
તો કિસ્સા જુદા પાડીએ:
\begin{enumerate}
  \item જો~$R$ એ
  મુખ્ય સૂત્ર~$!B \lif !C$
  ધરાવતું~\Intro{\lif}
  હોય અને~$x:!B$
  આધારવિધાનના સંદર્ભમાં
  હોય પણ નિષ્કર્ષના
  સંદર્ભમાં ન હોય,
  તો~$\delta$નો અંત
  આ રીતે થાય છે:
  \[
  \AxiomC{}
  \RightLabel{$\delta_1$}
  \Deduce$x:!B, \Gamma \fCenter !C$
  \RightLabel{\Intro{\lif}}
  \UnaryInf$\Gamma \fCenter !B \lif !C$
  \DisplayProof
  \]
  આગમન-પરિકલ્પનાથી
  $!B, \Gamma' \Sequent !C$ની
  (અથવા~$!C \ident \lfalse$
  હોય તો~$!B, \Gamma'
  \Sequent \ $ની)
  !!a{proof}~$\pi_1$
  મળે છે. પછી
  \RightR{\lif} લગાવી
  $\pi$ મળે છે
  ($!C \ident \lfalse$
  હોય તો પહેલાં
  \RightR{\Weakening}
  લગાવીએ).
  \item જો~$R$ એ
  \Intro{\lif} હોય,
  પણ~$x:!B$
  આધારવિધાનના સંદર્ભમાં
  ન હોય, તો~$\delta$નો
  અંત આ રીતે થાય છે:
  \[
  \AxiomC{}
  \RightLabel{$\delta_1$}
  \Deduce$\Gamma \fCenter !C$
  \RightLabel{\Intro{\lif}}
  \UnaryInf$\Gamma \fCenter !B \lif !C$
  \DisplayProof
  \]
  આગમન-પરિકલ્પનાથી
  $\Gamma' \Sequent !C$ની
  !!a{proof}~$\pi_1$
  મળે છે. તેમાં
  પહેલાં~$!B$ સાથે
  \LeftR{\Weakening}
  અને પછી~\RightR{\lif}
  લગાવી~$\pi$
  મેળવીએ છીએ.
  \item જો~$R$ એ
  \Elim{\lif} હોય,
  તો~$\delta$નો
  અંત આ રીતે થાય છે:
  \[
  \AxiomC{}
  \RightLabel{$\delta_1$}
  \Deduce$\Gamma_1 \fCenter !B \lif !A$
  \AxiomC{}
  \RightLabel{$\delta_2$}
  \Deduce$\Gamma_2 \fCenter !B$
  \RightLabel{\Elim{\lif}}
  \BinaryInf$\Gamma_1, \Gamma_2 \fCenter !A$
  \DisplayProof
  \]
  અહીં~$\Gamma = \Gamma_1 \cup \Gamma_2$.
  આગમન-પરિકલ્પનાથી
  $\Gamma_1' \Sequent !B \lif !A$ની
  !!{proof}~$\pi_1$
  અને~$\Gamma_2' \Sequent !B$ની
  સાબિતી~$\pi_2$
  મળે છે.
  !!{proof}~$\pi$
  આ રીતે મળે છે:
  \[
  \AxiomC{}
  \RightLabel{$\pi_1$}
  \Deduce$\Gamma_1' \fCenter !B \lif !A$
  \AxiomC{}
  \RightLabel{$\pi_2$}
  \Deduce$\Gamma_2' \fCenter !B$
  \Axiom$!A \fCenter !A$
  \RightLabel{\LeftR{\lif}}
  \BinaryInf$!B \lif !A, \Gamma_2' \fCenter !A$
  \RightLabel{\Cut}
  \BinaryInf$\Gamma_1', \Gamma_2' \fCenter !A$
  \DisplayProof
  \]
  બહુગણ~$\Gamma_1' \cup \Gamma_2'$
  કદાચ~$\Gamma' =
  (\Gamma_1 \cup \Gamma_2)'$
  જેટલો ન હોય:
  દાખલા તરીકે,
  $\Gamma_1$ અને~$\Gamma_2$
  બંનેમાં~$x:!C$ હોય,
  તો~$\Gamma_1' \cup \Gamma_2'$
  માં~$!C$ની બે નકલો
  છે, જ્યારે
  $(\Gamma_1 \cup \Gamma_2)'$
  માં એક જ છે.
  યોગ્ય
  \LeftR{\Contraction}
  અનુમાનો ઉમેરી
  આ ભેદ દૂર કરી શકાય.

  જો~$!A \ident \lfalse$
  હોય, તો~$\delta_2$ પર
  આગમન-પરિકલ્પના
  લગાવતાં
  $\Gamma_2' \Sequent !B$ની
  !!a{proof} મળે છે.
  તેને પહેલા આધારવિધાનની
  સાબિતી અને અસત્યના
  ખાલી-ઉત્તરાંગ સ્વયંસિદ્ધ
  સાથે~\LeftR{\lif}
  અને પછી~\Cut વડે
  જોડીને
  !!{proof}~$\pi$
  મળે છે; અહીં~$!A$
  માટે ઉત્તરાંગ ખાલી છે.
  \sourcecorrection{OLMD-250}{મૂળમાં અસત્ય-નિષ્કર્ષના કિસ્સામાં $\delta_2$નું અનુવાદિત ઉત્તરાંગ ખાલી કહેવાયું છે; તે હજી $!B$ સાબિત કરે છે. ખાલી ઉત્તરાંગ અસત્ય-ડાબા સ્વયંસિદ્ધમાંથી આવે છે, જેને અનુસરણ-ડાબા નિયમ અને કટથી જોડવું પડે.}
\item જો~$R$ એ
\FalseCl હોય,
તો~$\delta$નો
અંત આ રીતે થાય છે:
  \[
  \AxiomC{}
  \RightLabel{$\delta_1$}
  \Deduce$x: \lnot !A, \Gamma \fCenter \lfalse$
  \RightLabel{\FalseCl}
  \UnaryInf$\Gamma \fCenter !A$
  \DisplayProof
  \]
  આગમન-પરિકલ્પનાથી
  $\lnot !A,
  \Gamma' \Sequent \ $ની
  !!a{proof}~$\pi_1$
  છે. !!{proof}~$\pi$
  આ રીતે લઈએ:
  \[
  \Axiom$!A \fCenter !A$
  \RightLabel{\RightR{\lnot}}
  \UnaryInf$\fCenter !A, \lnot !A$
  \AxiomC{}
  \RightLabel{$\pi_1$}
  \Deduce$\lnot !A, \Gamma' \fCenter $
  \RightLabel{\Cut}
  \BinaryInf$\Gamma' \fCenter !A$
  \DisplayProof
  \]
  \sourcecorrection{OLMD-251}{મૂળના અંતિમ G2 વૃક્ષમાં $\pi_1$ના સિક્વન્ટમાં N2નું $x:$ ચિહ્ન રહી ગયું છે; G2 પૂર્વાંગ ચિહ્નવિહોણા સૂત્રોનું બહુગણ છે.}
\end{enumerate}
નોંધો કે ફક્ત
\FalseCl{}ના રૂપાંતરમાં
જમણી બાજુ એકથી વધુ
!!{formula} ધરાવતી
!!a{proof}~$\pi$
મળે છે. એટલે
\Log{N2i}ની
!!{proof}નું રૂપાંતર
\Log{G2i}-!!{proof}
છે.
\end{proof}

\end{document}
